\documentclass[11pt]{article}
\usepackage{amsmath,amssymb}

\title{A current release score for backtracking survey propagation}
\author{}
\date{}

\begin{document}
\maketitle

\begin{abstract}
The backtracking survey propagation algorithm of Marino, Parisi and Ricci-Tersenghi assigns variables so as to keep the complexity large, and it may later release an assignment. In the usual implementation the score used to choose that release is the bias stored at the time of the assignment. We replace that stored bias by a score computed from the messages of the current survey-propagation fixed point. The new score is the fraction of clusters still compatible with the assigned value. On 120 direct checks it tracks the change in complexity after a real release, with correlation $0.997$ and the correct sign in every case. At fixed clause density it makes the fall of the complexity slower and solves slightly more formulae than the stored-bias rule. A slow fall is not enough, by itself, to move the algorithmic threshold: at a higher density every rule we tried still failed.
\end{abstract}

\section{The 2016 algorithm}

We use the notation of Marino, Parisi and Ricci-Tersenghi \cite{marino2016}. A random $K$-SAT formula has $N$ Boolean variables and $M=\alpha N$ clauses of length $K$. The factor graph has a clause node $a$ for each clause and a variable node $i$ for each variable $x_i\in\{+1,-1\}$. The survey-propagation equations on the edges of this graph are
\begin{align}
\hat m_{a\to i}
&=
\prod_{j\in\partial a\setminus i} m_{j\to a},
\label{eq:sp1}
\\
m_{i\to a}
&=
\frac{\pi^-_{ia}\,(1-\pi^+_{ia})}{1-\pi^+_{ia}\,\pi^-_{ia}},
\label{eq:sp2}
\\
\pi^{\pm}_{ia}
&=
1-\prod_{b\in\partial^{\pm}_i a}(1-\hat m_{b\to i}).
\label{eq:sp3}
\end{align}
Here $\partial a$ is the set of variables in clause $a$. The set $\partial^+_i a$ (respectively $\partial^-_i a$) contains the clauses, other than $a$, that are satisfied (respectively violated) when $x_i$ is set to satisfy $a$. The message $\hat m_{a\to i}$ is the fraction of clusters in which clause $a$ forces $x_i$. The message $m_{i\to a}$ is the fraction of clusters in which $x_i$ is frozen to a value that does not satisfy $a$.

A fixed point $\{m^*_{i\to a},\hat m^*_{a\to i}\}$ gives the complexity, which is the logarithm of the number of clusters,
\begin{equation}
\Sigma
=
\sum_i \Sigma_i
+
\sum_a (1-K_a)\,\Sigma_a,
\label{eq:sigma}
\end{equation}
with $K_a$ the present length of clause $a$ and
\begin{align}
\Sigma_a
&=
\log\Bigl(1-\prod_{j\in\partial a} m^*_{j\to a}\Bigr),
&
\Sigma_i
&=
\log\bigl(1-\pi^+_i\pi^-_i\bigr),
\label{eq:sigma-parts}
\\
\pi^{\pm}_i
&=
1-\prod_{b\in\partial^{\pm}_i}(1-\hat m^*_{b\to i}).
\label{eq:pi}
\end{align}
The same fixed point gives the fraction of clusters in which $x_i$ is frozen positive, frozen negative, or not frozen:
\begin{equation}
w^{\pm}_i
=
\frac{\pi^{\pm}_i\,(1-\pi^{\mp}_i)}{1-\pi^+_i\pi^-_i},
\qquad
w^0_i
=
1-w^+_i-w^-_i.
\label{eq:w}
\end{equation}
Survey-inspired decimation assigns the more probable value, $x_i=+1$ if $w^+_i>w^-_i$ and $x_i=-1$ otherwise. That assignment multiplies the number of clusters by the bias
\begin{equation}
b_i
=
1-\min(w^+_i,w^-_i)
=
\max(w^+_i,w^-_i)+w^0_i.
\label{eq:bias}
\end{equation}
The corresponding change of complexity is $\Delta\Sigma\simeq\log b_i\le 0$. Variables with the largest $b_i$ are assigned, a fraction $f$ of them at each step. In the 2016 experiments $f=10^{-3}$. After each step the surveys are computed again on the residual formula.

Backtracking survey propagation adds the opposite move. A variable that has already been assigned may be released and assigned again later. The ratio of release steps to assignment steps is a constant $r\in[0,1)$. The choice $r=0$ is survey-inspired decimation. The running time grows as $1/(1-r)$ relative to decimation without releases, and it stays of order $N\log N$ for any fixed $r<1$. The 2016 rule uses the same bias for both moves: assign a free variable of large $b_i$, and release an assigned variable of small $b_i$. The paper states that this bias can be recomputed for an assigned variable from the surveys that arrive at it. If that bias, large at the assignment, later becomes small, the assignment is a candidate for release.

The algorithm stops in one of two ways. If $\Sigma$ becomes too small, or negative, the iteration may cease to converge, or the residual formula may contain an empty clause. If instead a trivial fixed point is reached, with every message equal to zero, the residual formula is given to WalkSAT. Solutions found in this way have no frozen variables under the whitening procedure.

\section{What the stored bias measures}

In the solver that accompanies the 2016 algorithm, the surveys $w^{\pm}_i$ and the bias $b_i$ are recomputed only for free variables. An assigned variable keeps the bias it had when it was fixed. A later release step sorts the assigned variables by those stored numbers.

Those numbers were computed on different residual formulae. A variable fixed early keeps a bias from a formula that still contained almost every clause. A variable fixed late keeps a bias from a much smaller formula. The comparison therefore mixes the present compatibility of the assignment with the age of the assignment. The rule written in the paper asks for the present bias. The stored value is only a proxy for it.

\section{The release score}

Let $u$ be the current partial assignment, and let $\varepsilon_k\in\{+1,-1\}$ be the value already given to an assigned variable $x_k$. Write $\mathcal N(u)$ for the number of clusters compatible with $u$, so that $\Sigma(u)=\log\mathcal N(u)$ at the survey-propagation estimate. Releasing $k$ produces a partial assignment $u\setminus k$ that allows both values of $x_k$. If a fraction $I_k$ of the clusters of $u\setminus k$ have $x_k=\varepsilon_k$, then
\begin{equation}
\mathcal N(u)
\simeq
I_k\,\mathcal N(u\setminus k),
\label{eq:Ik-def}
\end{equation}
and the gain on release is
\begin{equation}
\Delta\Sigma_{\mathrm{release}}(k)
=
\Sigma(u\setminus k)-\Sigma(u)
\simeq
-\log I_k.
\label{eq:dS}
\end{equation}
A small $I_k$ means that the assigned value is compatible with few of the clusters that the current surveys describe. Releasing it multiplies the number of clusters by about $1/I_k$.

The fraction $I_k$ is not the bias $b_k$. The two Boolean values do not form a probability on the clusters, because an unfrozen cluster is counted on both sides:
\begin{equation}
\bigl(w^+_k+w^0_k\bigr)+\bigl(w^-_k+w^0_k\bigr)
=
1+w^0_k.
\label{eq:not-prob}
\end{equation}
The fraction compatible with the value that was actually assigned is
\begin{equation}
I_k
=
\begin{cases}
1-w^-_k = w^+_k+w^0_k & \text{if }\varepsilon_k=+1,\\
1-w^+_k = w^-_k+w^0_k & \text{if }\varepsilon_k=-1.
\end{cases}
\label{eq:I}
\end{equation}
If the assigned value is still the majority value, then $I_k=b_k$. If the current surveys prefer the opposite value, then $I_k=1-\max(w^+_k,w^-_k)$, which is smaller than $b_k$. Sorting by $b_k$ releases the variables that look uncertain. Sorting by $I_k$ releases the variables whose assigned value the current surveys no longer support.

The weights $w^{\pm}_k$ in \eqref{eq:I} are not the weights stored at assignment time. They are computed from the current fixed point by \eqref{eq:pi} and \eqref{eq:w}, using the surveys that the neighbouring clauses send to $k$. The variable is not released in order to obtain these weights, and survey propagation is not run again. Only the ranking of the release step changes. The ratio $r$, the fraction $f$, the choice of which free variable to assign, and the choice of sign all stay as in the 2016 algorithm. We call this rule dynamic-$I$ backtracking.

Along any path the sum of the changes of $\Sigma$ depends only on the endpoints,
\begin{equation}
\sum_t \Delta\Sigma_t
=
\Sigma_{\mathrm{final}}-\Sigma_{\mathrm{initial}}.
\label{eq:telescope}
\end{equation}
Two paths with the same endpoints therefore have the same total loss. The quantity that can differ is the profile $\Sigma(d)$, where $d$ is the number of assigned variables. The aim of \eqref{eq:dS} is to avoid early, large drops of that profile. The desired end of a successful run is unchanged: $\Sigma$ stays positive while the surveys are nontrivial, and $\Sigma$ reaches zero together with the trivial fixed point. The failure mode is a drop of $\Sigma$ to zero, or below zero, while the surveys are still nontrivial, followed by loss of convergence or by an empty clause.

\section{Check against a real release}

Equation \eqref{eq:dS} was compared with the complexity obtained by actually releasing the variable. For each trial the variable was released, survey propagation was run to convergence on the enlarged formula, and $\Delta\Sigma$ was read from \eqref{eq:sigma}. This was done for 120 assigned variables.

The predicted gain $-\log I_k$ and the measured gain have correlation $0.99684$. The sign agrees in all 120 trials. The mean absolute error is $0.00056$. For these trials the local quantity $I_k$, built from the current incoming surveys, is a faithful estimate of the change in the global complexity. This is the strongest evidence we have for the score. It does not yet say that the solver will reach a higher clause density.

\section{Solver runs}

The release rule was then used inside the full algorithm, with the same $r$ as the usual backtracking runs. Every formula reported as solved was checked clause by clause. The sample sizes are small, and they are stated below so that the counts are not read as a measurement of the algorithmic threshold $\alpha_a$.

For $K=3$, $N=10^4$ and $\alpha=4.2$, five formulae were drawn. Dynamic-$I$ solved 5 out of 5. The stored certainty bias solved 4 out of 5. Polarization, $|w^+_i-w^-_i|$, solved 5 out of 5. Relative to certainty, the mean descent of $\Sigma$ was 26.0 percent smaller, the roughness of $\Sigma$ at a given depth was 17.7 percent smaller, and the largest drop at a given depth was 17.9 percent smaller. The running time was 5.8 percent larger.

For $K=4$, $N=10^3$ and $\alpha=9.5$, ten formulae were drawn. Dynamic-$I$ solved 6 out of 10. Certainty solved 5 out of 10. Polarization solved 6 out of 10. Relative to certainty, the mean absolute change of $\Sigma$ was 8.6 percent smaller, the mean descent was 23.2 percent smaller, the largest drop at a given depth was 39.8 percent smaller, and the roughness was 4.8 percent smaller. The running time was larger by a factor $1.94$. The eleven solved formulae from the dynamic-$I$ runs, five at $K=3$ and six at $K=4$, all passed the clause check.

For $K=4$ and $\alpha=9.7$, five further formulae were drawn. Dynamic-$I$ still produced the smoothest complexity curve. Every method failed on all five formulae, including dynamic-$I$.

\begin{table}[ht]
\centering
\begin{tabular}{lccc}
\hline
setting & dynamic-$I$ & certainty $b_i$ & polarization \\
\hline
$K=3$, $N=10^4$, $\alpha=4.2$, 5 seeds & 5/5 & 4/5 & 5/5 \\
$K=4$, $N=10^3$, $\alpha=9.5$, 10 seeds & 6/10 & 5/10 & 6/10 \\
$K=4$, $N=10^3$, $\alpha=9.7$, 5 seeds & 0/5 & 0/5 & 0/5 \\
\hline
\end{tabular}
\caption{Number of formulae solved. The release rule is the only change in the dynamic-$I$ column. Sample sizes are the number of seeds in the first column.}
\end{table}

\section{Changes that did not help}

Several other uses of $\Sigma$ were tried and did not improve on the 2016 ranking.

A one-step lookahead assigns, in a copy of the formula, each of the few free variables with the largest $b_i$, runs survey propagation again, and keeps the assignment with the highest residual $\Sigma$. This is a direct attempt to minimize the immediate drop. On formulae with a few hundred variables the worst nontrivial slope of $\Sigma$ changed only in the third digit, while the number of survey iterations grew by a factor of about three. A $K=4$ run that lost convergence under the 2016 rule also lost convergence under the lookahead. Equation \eqref{eq:telescope} says why a local reward on $\Delta\Sigma$ is a weak guide: paths that meet at the same end have the same total loss. The lookahead also trusts $\Sigma$ in the region where the estimate is least reliable. Survey propagation can converge with $\Sigma>0$ on a formula that has no solution, and it can fail to converge on a formula that still has solutions.

Hand-made scores built from $w^{\pm}_i$ and $w^0_i$, including polarization and a one-parameter family between polarization and certainty, gave algorithmic densities that agreed with certainty within the resolution of those scans. Forbidding two variables in the same clause from being fixed in one batch did not change the outcome at the sizes where the batch contains one variable, which is the usual situation for $N\lesssim 10^3$ at $f\simeq 10^{-3}$.

Raising $r$ only after a steep drop of $\Sigma$, or after a slow convergence, also left the nontrivial part of the profile unchanged on the runs we did. During that part of the run the relative drop per step stayed near one or two percent, below the threshold that would have added releases. The extra releases appeared only after $\Sigma$ had already collapsed.

An exact satisfiability solver was used, on small formulae, to flip the sign of an assignment when the survey sign made the residual formula unsatisfiable. That correction removed some fatal assignments. The runs that remained unsolved then failed because survey propagation did not converge, even though the residual formula was still satisfiable. Choosing a better sign does not repair a fixed point that the iteration cannot reach.

\section{What can be claimed}

The 2016 paper already asks for a present bias of each assigned variable, and it already uses backtracking to undo an assignment whose bias has decayed. The quantity $I_k$ is the form of that idea that matches the change in the number of clusters, because it uses the assigned sign and the unfrozen weight, as in \eqref{eq:I}. The measurements say three things.

First, $-\log I_k$ computed from the current incoming surveys agrees with the complexity change of a real release, at the precision given in the check above.

Second, using $I_k$ inside the solver slows the fall of $\Sigma(d)$ and, at the two densities where some formulae are solved, solves slightly more formulae than the stored bias. The time overhead is small at $K=3$ and about a factor of two at $K=4$.

Third, a smoother profile is not a proof that the algorithmic threshold has moved. At $\alpha=9.7$ for $K=4$ the smoothest profile still produced no solution. Keeping many clusters alive is favourable. The solver also has to remain on a survey fixed point that converges, and on clusters that the final local search can finish. The complexity at $m=0$ counts clusters equally, and the numerous clusters are not always the ones that can be completed in linear time. The present evidence supports a better release rule at fixed $\alpha$. It does not support a new value of $\alpha_a$.

A claim about the threshold would need a sweep in $\alpha$, several sizes $N$, and enough seeds to attach an uncertainty to the density at which the success rate falls. Those runs have not been done.

\begin{thebibliography}{9}
\bibitem{marino2016}
R.~Marino, G.~Parisi and F.~Ricci-Tersenghi,
The backtracking survey propagation algorithm for solving random $K$-SAT problems,
Nature Communications \textbf{7}, 12996 (2016).
\bibitem{parisi2003}
G.~Parisi,
A backtracking survey propagation algorithm for $K$-satisfiability,
arXiv:cond-mat/0308510 (2003).
\bibitem{mezard2002}
M.~M\'ezard, G.~Parisi and R.~Zecchina,
Analytic and algorithmic solution of random satisfiability problems,
Science \textbf{297}, 812 (2002).
\end{thebibliography}

\end{document}
